\section{Neural and Statistical Baselines}

The baselines in this section compare the template methods with two learned classifiers: a feature-based MLP and a small CNN. Both models train on the training split, use the development split for thresholding or early stopping, and report final metrics on the test split. Although the task is one-hot four-way classification, both models keep the same four-output interface as the template methods and use $\arg\max$ for the final label.

\subsection{MLP on MFCC and band-energy statistics}

The MLP does not take a full spectrogram as input. It uses a fixed vector of summary features. For each recording, the pipeline computes log-mel features, extracts 13 MFCC coefficients and their first-order differences, and concatenates
\[
z_i=
\left[
\mathrm{mean}(\mathrm{MFCC}),
\mathrm{std}(\mathrm{MFCC}),
\mathrm{mean}(\Delta\mathrm{MFCC}),
\mathrm{std}(\Delta\mathrm{MFCC}),
\mathrm{bandEnergy}
\right].
\]
The first four groups contribute $13\times4=52$ dimensions. The band-energy block contributes 10 more, so $z_i\in\mathbb{R}^{62}$. The features are standardized with the training-set mean and variance.

The network has two hidden layers:
\[
h_1=\mathrm{ReLU}(W_1z_i+b_1),\qquad
h_2=\mathrm{ReLU}(W_2h_1+b_2),\qquad
p_i=\sigma(W_3h_2+b_3),
\]
with widths 128 and 64. The loss is binary cross entropy over the four output labels plus $L_2$ regularization:
\[
\mathcal{L}_{\mathrm{MLP}}
=-\sum_{i,c}\left[y_{i,c}\log p_{i,c}+(1-y_{i,c})\log(1-p_{i,c})\right]
+\lambda\|\theta\|_2^2.
\]
The MLP tests whether whole-word MFCC statistics already contain enough class information. It does not localize the initial consonant, so vowel and speaker differences remain mixed with the target cue. Appendix~\ref{alg:mlp-ai} gives the algorithm.

\subsection{CNN on whole-word log-mel spectrograms}

The CNN receives a whole-word log-mel spectrogram $L_i\in\mathbb{R}^{T\times40}$. The input is standardized by the training-set global mean and standard deviation and then passed as a single-channel image. The architecture is
\[
\mathrm{Conv}(1,24)\rightarrow \mathrm{BN}\rightarrow \mathrm{ReLU}\rightarrow \mathrm{MaxPool},
\]
\[
\mathrm{Conv}(24,48)\rightarrow \mathrm{BN}\rightarrow \mathrm{ReLU}\rightarrow \mathrm{MaxPool},
\]
\[
\mathrm{Conv}(48,96)\rightarrow \mathrm{BN}\rightarrow \mathrm{ReLU}\rightarrow
\mathrm{GlobalAvgPool}\rightarrow \mathrm{FC}(96)\rightarrow \mathrm{Dropout}(0.25)\rightarrow \mathrm{FC}(4).
\]
Let the logits be $u_{i,c}$ and probabilities be $p_{i,c}=\sigma(u_{i,c})$. The loss is weighted binary cross entropy:
\[
\mathcal{L}_{\mathrm{CNN}}
=-\sum_{i,c}\alpha_c y_{i,c}\log p_{i,c}
-(1-y_{i,c})\log(1-p_{i,c}),
\]
where $\alpha_c$ is computed from the positive/negative ratio in the training split and clipped to $[1,10]$. Training uses AdamW with learning rate $10^{-3}$, weight decay $10^{-4}$, at most 10 epochs, and development macro-F1 for early stopping.

The CNN lets the model learn local patterns from the whole spectrogram. This assumption needs more data than the 3000-utterance, 60-word corpus provides, and the model is not forced to attend to the word onset.

\begin{table}[H]
    \centering
    \caption{Overall metrics for the learned baselines.}
    \label{tab:ai-summary}
    \begin{tabular}{lccccc}
        \toprule
        Method & Label acc. & Exact match & Macro-F1 & Micro-F1 & Samples-F1 \\
        \midrule
        MLP & 0.7633 & 0.5267 & 0.5258 & 0.5267 & 0.5267 \\
        CNN & 0.6600 & 0.3200 & 0.2219 & 0.3200 & 0.3200 \\
        \bottomrule
    \end{tabular}
\end{table}

The MLP is a useful statistical baseline, but it remains below the extended template method. The CNN trains successfully but performs worse than the strict FFT template on exact-match accuracy. In this dataset, a small end-to-end spectrogram model does not recover the onset cue more reliably than a hand-designed local spectral comparison.

\begin{figure}[H]
    \centering
    \includegraphics[width=0.92\textwidth]{figures/mswc_course_gbdz_local_patch_v1/cnn_history.png}
    \caption{CNN training history.}
    \label{fig:cnn-history}
\end{figure}

\begin{figure}[H]
    \centering
    \includegraphics[width=0.92\textwidth]{figures/mswc_course_gbdz_local_patch_v1/metric_summary.png}
    \caption{Overall metric summary for all methods.}
    \label{fig:metric-summary}
\end{figure}

\begin{figure}[H]
    \centering
    \includegraphics[width=0.92\textwidth]{figures/mswc_course_gbdz_local_patch_v1/per_label_f1.png}
    \caption{Per-label F1 comparison across all methods.}
    \label{fig:per-label-f1}
\end{figure}
